\documentclass[10pt,a4paper]{article}
\usepackage[margin=1.25cm,top=1.0cm,bottom=0.85cm]{geometry}
\usepackage[T1]{fontenc}
\usepackage{mathptmx}
\usepackage{xcolor,hyperref,enumitem,tabularx}
\pagestyle{empty}
\definecolor{navy}{RGB}{26,62,102}
\hypersetup{colorlinks=true,urlcolor=blue,linkcolor=blue,pdfborder={0 0 0}}
\setlength{\parindent}{0pt}
\setlength{\parskip}{0pt}
\setlist[itemize]{leftmargin=14pt,itemsep=0pt,topsep=2pt,parsep=0pt}
\newcommand{\sectiontitle}[1]{\par\vspace{3pt}\textcolor{navy}{\large\textbf{#1}}\par\vspace{-3pt}\textcolor{navy}{\rule{\linewidth}{0.45pt}}\vspace{0pt}}
\newcommand{\entry}[2]{\textbf{#1}\hfill{\color{black!65}#2}\par}
\begin{document}\fontsize{9.5}{10.5}\selectfont
\begin{center}
{\fontsize{19}{20}\selectfont\textcolor{navy}{\textbf{Xuan Li Feng}}}\par\vspace{2pt}
\href{mailto:xl4723@ic.ac.uk}{xl4723@ic.ac.uk}\quad | \quad London, United Kingdom\quad | \quad MEng Biomedical Engineering, Imperial College London\par
\vspace{1pt}\href{https://xuan-lifeng.github.io}{Research portfolio: xuan-lifeng.github.io}
\end{center}\vspace{-5pt}
\sectiontitle{EDUCATION}
\entry{MEng Biomedical Engineering}{2023--2027 (expected)}
\textit{Imperial College London}
\begin{itemize}
\item Expected graduation: June 2027.
\item Relevant coursework: Image Processing; Digital Biosignal Processing; Statistical Signal Processing \& Inference; Biomedical Instrumentation.
\end{itemize}
\sectiontitle{RESEARCH EXPERIENCE}
\entry{MEng Thesis: Anatomy-Aware Surgical Robotics}{2026--2027}
\textit{Imperial College London \quad | \quad Supervised by Dr Stuart Bowyer}
\begin{itemize}
\item Developing a stereoscopic near-infrared fluorescence (NIRF) pipeline to reconstruct 3D vascular geometry from paired images for anatomy-aware surgical assistance.
\item Current work focuses on vessel segmentation, stereo correspondence, triangulation, depth estimation, centreline extraction, and validation against synthetic ground truth and tissue-mimicking phantoms.
\item Evaluating reconstruction reliability using depth and geometry error metrics; investigating failure modes from segmentation and correspondence uncertainty.
\end{itemize}
\entry{Deep Neural Receptive Fields in Mouse Auditory Cortex}{2025--2026}
\textit{Imperial College London \quad | \quad Supervised by Dr A. S. Kozlov}
\begin{itemize}
\item Designed and trained a five-layer causal Conv1D model to predict auditory-cortical spiking responses from naturalistic ultrasonic vocalisations in mice.
\item Used Jacobian-based automatic differentiation to derive dynamic spectrotemporal receptive fields and compare neural-network behaviour with STA, STC, MNE, and dSTRF analyses.
\end{itemize}
\entry{Research Internship: Magnetoelastic Energy Harvesting}{Jun--Aug 2025}
\textit{Zhejiang University \quad | \quad Principal Investigator: Professor Jikui Luo}
\begin{itemize}
\item Conducted full-time research on GaIn-doped PVDF composites, using COMSOL multiphysics simulation and molecular-dynamics and density-functional-theory analysis to support experimental interpretation.
\end{itemize}
\sectiontitle{PROJECTS}
\entry{Wearable Eye-Tracking AAC Device}{2024--2025}
\textit{Imperial College London}
\begin{itemize}
\item Developed gaze and head-tilt input for a wearable AAC device using OpenCV, Raspberry Pi Zero 2 W, IMU sensing, and Google Glass; achieved 78.2\% gaze-selection accuracy and 91.8\% head-tilt accuracy with latency below 100 ms.
\end{itemize}
\entry{Remote Patient Monitoring System}{2025}
\textit{Imperial College London}
\begin{itemize}
\item Developed components of a Java Swing clinical monitoring system with REST API and PostgreSQL backend, including real-time waveform display and digital-twin visualisation.
\end{itemize}
\entry{Biomedical Instrumentation Systems}{2024--2026}
\textit{Imperial College London}
\begin{itemize}
\item Designed and validated multistage analogue cardiac and respiratory sensing systems using LM324 operational amplifiers, active filters, SPICE, and hardware prototyping; verified 3700 voltage gain and real-time monitoring at 30 breaths per minute.
\end{itemize}
\sectiontitle{PUBLICATION}
Lu, J., Zhang, K., \ldots, \textbf{Li Feng, X.}, \ldots et al. (2025). \textit{GaIn Induced Polarization Control in PVDF for High-Efficiency Energy Harvesting and Instantaneous Wireless Sensing.} Energy \& Environmental Materials. \href{https://doi.org/10.1002/eem2.70201}{doi:10.1002/eem2.70201}
\sectiontitle{TECHNICAL SKILLS}
\renewcommand{\arraystretch}{1.05}
\begin{tabularx}{\linewidth}{@{}>{\bfseries}p{3.05cm}X@{}}
Medical Imaging \& Vision & Stereo vision, 3D triangulation, image segmentation, centreline extraction, reconstruction validation, NIRF imaging concepts, OpenCV\\
Machine Learning \& Neural Data & PyTorch, Conv1D, automatic differentiation/Jacobians, cross-validation, neural response prediction, model evaluation\\
Signal Processing \& Statistics & Wiener and LMS filtering, AR modelling, spectral estimation, Cram\'er--Rao bounds, time-series analysis, hypothesis testing, Pearson correlation\\
Programming \& Research Tools & Python, MATLAB, Java, JavaScript, Kotlin, Git/GitHub, Jupyter, PostgreSQL\\
Hardware \& Systems & Embedded vision, Raspberry Pi, IMU sensing, PCB design, analogue filtering and amplification, SPICE, COMSOL Multiphysics
\end{tabularx}
\end{document}
